% Section 1: Introduction (~1,200 words)

\subsection{Why Validation Matters for Algorithmic Redistricting}

Congressional redistricting shapes the composition of the U.S. House of
Representatives for an entire decade. The stakes---who governs, whose
preferences are amplified, whose communities are intact---are high enough
that every methodological choice requires rigorous justification.
Algorithmic approaches to redistricting promise to replace the backroom
negotiations and partisan calculations of legislative map-drawing with
reproducible, transparent, geometric optimization. But a promise of
transparency is not the same as a demonstration of robustness.

Consider what could go wrong even for a technically sound algorithm.
A redistricting method that works well with 2020 Census tracts might
produce different results with 2010 or 2000 tract boundaries, which
differ in their shapes and sizes due to decennial boundary revisions.
It might produce very different districts depending on whether it uses
census blocks ($\sim$100 residents each), block groups ($\sim$1,200),
or tracts ($\sim$4,000)---a concern known in geographic information
science as the Modifiable Areal Unit Problem
\citep{openshaw1984modifiable}. The algorithm might behave well for
states with a long political history of non-partisan redistricting but
fail to generalize to states with extreme demographic concentration.
And even if none of these spatial or temporal failures occur, the
algorithm might---unintentionally---produce plans that systematically
favor one party over another, not through any defect in the code but
through the interaction of geometric optimization with political
geography \citep{chen2013unintentional}.

Validation is the discipline that answers these questions empirically
rather than by assertion. A validated algorithm is one that has been
tested against multiple failure modes and found robust. The C-track
(C.1 through C.5) of the Algorithmic Redistricting Research Portfolio
constitutes exactly such a validation program.

\subsection{The Five Validation Dimensions}

The C-track organizes validation along five dimensions, each addressed
by a dedicated empirical paper:

\begin{enumerate}[leftmargin=*]
  \item \textbf{Spatial robustness (C.1).} Does the algorithm produce
    consistent results when the spatial building blocks change? We test
    this by running the full redistricting pipeline at three resolutions:
    census tracts (baseline; $\approx$74,000 units nationally), block
    groups ($\approx$239,000 units), and census blocks ($\approx$8.1M
    units)---a 130$\times$ range in unit count
    \citep{dellalibera2026c1}.

  \item \textbf{Cross-census consistency (C.2).} Does the algorithm
    behave consistently across the three available census years
    (2000, 2010, 2020), even though tract boundaries, population
    distributions, and apportionment change between decades?
    C.2 introduces a slice-based validation framework that controls
    for geographic confounds \citep{dellalibera2026c2}.

  \item \textbf{Temporal stability (C.3).} When the same algorithm is
    applied to successive census years for the same state, how much do
    the resulting district boundaries change? Geographic stability---
    not just metric stability---is a legal and practical requirement
    \citep{dellalibera2026c3}.

  \item \textbf{Longitudinal analysis (C.4).} Over twenty years of
    census cycles, is the gap between algorithmic and enacted plans
    shrinking as redistricting reform spreads? C.4 places the algorithm
    in historical context, showing that the algorithmic approach becomes
    increasingly competitive as reform-era enacted maps improve
    \citep{dellalibera2026c4}.

  \item \textbf{Partisan fairness (C.5).} A neutral algorithm should
    produce plans with near-zero efficiency gap as a byproduct of
    geometric optimization alone. C.5 documents that algorithmic plans
    achieve a mean $|EG| = 0.04$, compared to $0.08$ for enacted maps,
    and does so without ever accessing partisan data
    \citep{dellalibera2026c5}.
\end{enumerate}

\subsection{The DIA Reproducibility Standard}

The Districting Integrity Act (DIA)\footnote{The Districting Integrity
Act is a model statute proposed in B.02 of this research portfolio
\citep{b02paper}; it has not been enacted by any legislature and is not
a binding legal requirement. References throughout this paper to DIA
``requirements'' describe the standards the model statute would impose
if adopted.} specifies that a legally defensible redistricting algorithm
must satisfy three reproducibility requirements: (R1) deterministic output from a
specified seed, (R2) invariance to input data perturbations within
defined tolerance bounds, and (R3) documented behavior across the
full range of geographies and census years for which the algorithm
is used. The C-track papers collectively provide the empirical
evidence for (R2) and (R3). B.7 \citep{b7paper} addresses (R1) by
characterizing seed sensitivity across 10,000 seeds per state.

\subsection{Paper Structure}

The remainder of this paper proceeds as follows.
Section~\ref{sec:spatial} summarizes C.1's MAUP findings.
Section~\ref{sec:cross-census} summarizes C.2's cross-census
validation framework and results.
Section~\ref{sec:temporal} synthesizes the temporal stability
findings from C.3 and the longitudinal trends from C.4.
Section~\ref{sec:partisan} presents the partisan fairness validation
from C.5.
Section~\ref{sec:synthesis} articulates the four properties that the
C-track collectively validates and connects them to DIA requirements.
Section~\ref{sec:conclusion} concludes.
